% Fichier complété par tools/complete_missing_corrections.py.
% Source : PPM/PPM-01/1020_PompeEnsieta/1020_PompeEnsieta.tex
% Statut : rédigé (7 question(s)).
\begin{corrigebox}[Corrigé rédigé]
\CorrigeQuestion{1}
Non. Une bille de même diamètre que l'alésage ne peut pas se déplacer librement et risque de se coincer. Il faut un jeu radial fonctionnel et un siège de diamètre inférieur à celui de la bille afin d'assurer l'étanchéité par contact annulaire.
\CorrigeQuestion{2}
Non. L'alésage recevant le raccord n'est qu'une surface de positionnement/étanchéité statique, alors que l'alésage du piston est une surface de guidage en mouvement. Ce dernier exige une rugosité plus faible, une meilleure cylindricité et une tolérance dimensionnelle plus serrée.
\CorrigeQuestion{3}
Il faut un ajustement avec jeu faible, de type glissant précis (par exemple H7/g6 ou H7/f7 selon la vitesse et la lubrification). Un serrage interdirait le mouvement et un ajustement trop lâche provoquerait fuite et désalignement.
\CorrigeQuestion{4}
Le composant prisonnier ne peut être introduit si son diamètre maximal est supérieur au passage disponible après réalisation monobloc du corps. Il faut prévoir un bouchon démontable, un alésage débouchant, une pièce rapportée ou modifier l'ordre de montage. Sur le dessin, c'est le composant du clapet fermé par le corps qui doit être rendu accessible.
\CorrigeQuestion{5}
Non. Du cote tige, la section hydraulique utile est annulaire \(S_a=S_p-S_t\), tandis que, du cote oppose, elle vaut \(S_p\). Pour une meme course, les volumes sont \(V_p=S_pL\) et \(V_a=(S_p-S_t)L\); leur difference vaut \(S_tL\).
\CorrigeQuestion{6}
Pour une course $L$, le volume refoulé à la sortie 4 est $V_4=S_{\nm utile}L$. Si la sortie est côté plein piston, $V_4=\pi D^2L/4$; si elle est côté tige, $V_4=\pi(D^2-d^2)L/4$. Le débit moyen vaut ce volume multiplié par la fréquence des courses utiles.
\CorrigeQuestion{7}
Le dessin de définition du corps reprend toutes les surfaces
        fonctionnelles du dessin d'ensemble : alésage de guidage du piston avec tolérance
        et rugosité serrées, logements des clapets et sièges coniques, taraudages des
        raccords, épaulements de montage et perçages de communication. Les axes sont
        cotés depuis un même système de références et les formes intérieures sont montrées
        en coupe avec hachures réglementaires.
\end{corrigebox}
